\section{CUDA GPU并行编程：使用WMMA API编程张量核心}

\subsection{引言}
本节文章将深入讲解专用于张量核心（Tensor Cores）编程的 \textbf{WMMA API}（WarpMatrix Multiply-Accumulate）。在正式介绍API之前，我们首先回顾GPU执行矩阵乘法的四个通用步骤，以便
理解WMMA如何映射到这些基础操作中。

\subsection{矩阵乘法的核心步骤回顾}
无论是使用普通CUDA核心还是张量核心，矩阵乘法内核均遵循以下四个阶段：
\begin{enumerate}
    \item \textbf{分块（Tiling/Partitioning）}：根据矩阵维度$N$和块大小（Block Size）计算网格与线程块的划分。
    \item \textbf{加载（Loading）}：将子矩阵A和B从全局内存加载至共享内存（Shared Memory），以减少访存延迟。注意：输出矩阵C通常不需要预加载，除非进行累加操作。
    \item \textbf{计算（Execution）}：执行乘加运算。在传统实现中，结果暂存于寄存器变量（如 \texttt{sum}）中；在WMMA中，则存储于寄存器碎片（Register Fragments）中。
    \item \textbf{存储（Storing）}：将寄存器中的最终结果写回全局内存中的矩阵C。
\end{enumerate}

\subsection{WMMA API 核心概念}
\subsection{定义与适用架构}
WMMA 全称为 \textbf{Warp Synchronous MatrixMultiply-Accumulate}。
与普通CUDA核心以线程块（Thread Block）为调度单位不同，WMMA是以 \textbf{Warp（线程束）} 为基本执行单元的矩阵运算接口。

\begin{itemize}
    \item \textbf{硬件要求}：仅适用于计算能力（Compute Capability）$\geq 7.0$ 的设备。
    \item \textbf{架构背景}：张量核心首次引入于2017年的 Volta 架构（CC 7.x）。后续的 Turing、Ampere、Hopper 等架构均支持并扩展了该特性。Volta 之前的架构因缺乏物理张量核心，无法使用此API。
\end{itemize}

\subsection{碎片（Fragments）与数据布局}
WMMA 引入了“碎片”（Fragment）的概念，即分布在 Warp 内32个线程寄存器中的矩阵子块。

\subsection{支持的维度与类型}
WMMA 不支持任意尺寸的分块，必须使用预定义的形状。常见组合包括：
\begin{itemize}
    \item $16 \times 16 \times 16$
    \item $32 \times 8 \times 16$ / $8 \times 32 \times 16$
    \item 数据类型组合需严格遵循官方文档表格（例如：输入FP16 + 累加器FP32，或输入FP16 + 累加器FP16）。
    \item 支持的新类型包括 BF16、TF32、INT8 等，其中 TF32 具有与 FP32 相同的动态范围但精度较低（10位尾数）。
\end{itemize}

\subsection{内存布局：行主序 vs 列主序}
\begin{itemize}
    \item \textbf{Row-Major}：按行连续存储。
    \item \textbf{Col-Major}：按列连续存储。
\end{itemize}
\textbf{优化提示}：为最大化内存合并访问效率，通常建议矩阵A使用 Col-Major（便于按列加载），矩阵B使用Row-Major（便于按行加载），具体取决于内核的遍历策略。错误的布局会导致非合并访问，显著增加缓存未命中率和内存请求次数。

\subsection{WMMA 关键指令}
\begin{lstlisting}
// 1. 声明碎片
wmma::fragment<wmma::matrix_a, M, N, K, half, wmma::col_major> frag_a;
wmma::fragment<wmma::matrix_b, M, N, K, half, wmma::row_major> frag_b;
wmma::fragment<wmma::accumulator, M, N, K, float> frag_c;

// 2. 初始化累加器
wmma::fill_fragment(frag_c, 0.0f);

// 3. 从全局/共享内存加载到寄存器碎片
wmma::load_matrix_sync(frag_a, ptr_a, lda);
wmma::load_matrix_sync(frag_b, ptr_b, ldb);

// 4. 执行矩阵乘加: D = A * B + Cwmma::mma_sync(frag_c, frag_a, frag_b, frag_c);

// 5. 将结果存回内存
wmma::store_matrix_sync(ptr_c, frag_c, ldc, wmma::mem_row_major);
\end{lstlisting}

\subsection{Warp ID 的计算方法}
CUDA 提供了 \texttt{threadIdx} 和\texttt{blockIdx}，但\textbf{没有}预
定义的 \texttt{warpId} 变量。在使用WMMA 时，必须手动计算当前线程所属的 Warp ID。

\subsection{一维网格/线程块}
若线程块为一维组织，Warp ID 计算公式为：
\begin{equation}
    \text{warpId} = \frac{\text{threadIdx.x}}{32}
+ \left(\frac{\text{blockDim.x}}{32} \right) \times\text{blockIdx.x}
\end{equation}
\textbf{解析}：
\begin{itemize}
    \item $\text{threadIdx.x} / 32$：计算块内 Warp 偏移（整数除法）。
    \item $(\text{blockDim.x} / 32) \times \text{blockIdx.x}$：计算当前块之前所有块包含的 Warp 总数。
\end{itemize}

\subsection{二维网格扩展}
对于二维网格，需先将二维 blockIdx 展平为一维线性索引：
\begin{equation}
    \text{linearBlockId} = \text{blockIdx.y} \times\text{gridDim.x} + \text{blockIdx.x}
\end{equation}
随后代入上述公式即可。注意：Warp 的划分始终基于线程的线性索引，与网格维度无关。

\subsection{WMMA 矩阵乘法执行流程}
结合 Warp ID 计算，完整的 WMMA GEMM 内核逻辑如下：

\begin{enumerate}
    \item \textbf{计算 Warp 坐标}：根据 \texttt{warpId} 确定当前 Warp 负责处理的全局矩阵分块位置（Row, Col）。
    \item \textbf{初始化}：使用 \texttt{fill\_fragment} 将累加器碎片置零。
    \item \textbf{K维循环}：沿 K 维度以步长16（或其他碎片K维度）遍历：
    \begin{itemize}
        \item 计算子矩阵 A 和 B 在当前 K 步的指针偏移。偏移量依赖于 \texttt{Leading Dimension (LDA/LDB)} 而非简单的矩阵宽度 N。
        \item 使用 \texttt{load\_matrix\_sync} 加载数据到寄存器碎片。
        \item 调用 \texttt{mma\_sync} 执行乘加运算。
    \end{itemize}
    \item \textbf{写回结果}：循环结束后，使用 \texttt{store\_matrix\_sync} 将累加器碎片写入全局内存矩阵 C。
\end{enumerate}

\subsection{重要提示}
\begin{itemize}
    \item \textbf{Leading Dimension (LD)}：在计算内存偏移时，务必使用 LD 参数（通常为矩阵的实际分配行数/列数），而非逻辑维度。这是避免内存越界和数据错位的关键。
    \item \textbf{文档查阅}：WMMA 对数据类型、布局和尺寸的组合有严格限制。开发时请务必查阅 CUDA Programming Guide 中的 "Warp Matrix Functions" 章节，确认所用组合受硬件支持。
    \item \textbf{性能调优}：WMMA 仅是基础原语。实际高性能实现还需结合共享内存分块、向量化加载、流水线重叠等优化技术。
\end{itemize}

